\documentclass[11pt]{article}
\usepackage[margin=1in]{geometry}
\usepackage{amsmath,amssymb}
\usepackage{booktabs}
\usepackage{hyperref}
\usepackage{xcolor}
\title{Replication Report: \textit{Kondo Destruction and Multipolar Order}\\
\large arXiv:1807.09258 (Lai, Nica, Hu, Gong, Paschen, Si, 2018)}
\author{Automated replication (8-artifact standard)}
\date{2026-07-19}
\begin{document}
\maketitle

\section{Paper summary}
Lai \emph{et al.}\ study a spin-1 Kondo lattice model on the square lattice with
\emph{both} spin (dipole) and quadrupole local degrees of freedom,
\begin{align}
H_{\rm KL} &= H_S + H_c + H_K,\\
H_S &= \sum_{ij}\Big[ J_{ij}(\mathbf S_i\!\cdot\!\mathbf S_j)
          + K_{ij}(\mathbf S_i\!\cdot\!\mathbf S_j)^2\Big],\\
H_K &= \sum_j\big(J_K^{I}\,\mathbf S_j\!\cdot\!\mathbf s_{c,j}
          + J_K^{II}\,\mathbf Q_j\!\cdot\!\mathbf q_{c,j}\big).
\end{align}
Using DMRG they establish a $(\pi,\pi)$ \emph{antiferroquadrupolar} (AFQ) phase for
$J_1=1,\,K_1=1.2,\,J_2=0,\,K_2<0$, characterized by a staggered
$\langle Q_j^{x^2-y^2}\rangle\sim(-1)^j$. They then derive a quantum non-linear
sigma model (NL$\sigma$M) for the AFQ order using an SU(3) time-reversal-invariant
director ($d$-vector) representation and a combined boson--fermion renormalization
group (RG) analysis, concluding that the Kondo couplings are \emph{exactly marginal}
in the AFQ phase (Kondo destruction): the Kondo vertex carries positive powers of
$\sqrt{\Lambda/K_F}$ and is asymptotically suppressed.

The ``multipolar texture'' here is the spatial pattern of complex directors
$\mathbf d_j=(d_j^x,d_j^y,d_j^z)$ encoding the quadrupolar order:
$\mathbf d^{\rm gs}_A=(1,0,0)^T,\ \mathbf d^{\rm gs}_B=(0,1,0)^T$ on the two
sublattices (Fig.~2a of the paper).

\section{Scope of this replication}
Full DMRG (4000 SU(2) states, $L_y=8$ cylinders) and the field-theoretic RG are out
of scope for a laptop. We replicate the \emph{exactly checkable algebraic core} plus a
\emph{tractable finite-cluster ED surrogate} of the DMRG structure-factor claim:
\begin{itemize}
\item the operator algebra (biquadratic identities, quadrupole operators),
\item the $d$-vector $\to$ spin/quadrupole maps (Eqs.~S16--S22),
\item the ground-state director pattern and its energetics (Eq.~S23/S24),
\item the SU(3) Gell-Mann rotation structure and the flat (marginal) direction,
\item small-cluster ED of the spin-1 BLBQ model showing the $(\pi,\pi)$ AFQ peak.
\end{itemize}

\section{Checked claims and results}

\subsection*{C1 --- Biquadratic operator identities (exact)}
The paper's algebraic identities
\[(\mathbf S_i\!\cdot\!\mathbf S_j)^2=\tfrac12\mathbf Q_i\!\cdot\!\mathbf Q_j
-\tfrac12\mathbf S_i\!\cdot\!\mathbf S_j+\tfrac13\mathbf S_i^2\mathbf S_j^2,\qquad
\mathbf Q_i\!\cdot\!\mathbf Q_j = 2(\mathbf S_i\!\cdot\!\mathbf S_j)^2
+\mathbf S_i\!\cdot\!\mathbf S_j-\tfrac83\]
were verified operator-by-operator on the $9$-dimensional two-site spin-1 space:
residuals $\sim10^{-16}$--$10^{-15}$. \textbf{PASS.}

\subsection*{C2 --- $d$-vector maps (Eqs.~S16--S22)}
For $200$ random normalized directors we compared $\langle S^\alpha\rangle,
\langle Q^a\rangle$ from the $3\times3$ operators (in the TRI basis, Eq.~S10) against
the closed-form $d$-vector formulas $S^\alpha=-i\epsilon^{\alpha\beta\gamma}
\bar d^\beta d^\gamma$, $Q^{x^2-y^2}=-|d^x|^2+|d^y|^2$, etc. Max discrepancy
$5.6\times10^{-16}$. The ground-state directors carry \emph{zero} dipole moment.
\textbf{PASS.}

\subsection*{C3 --- $(\pi,\pi)$-AFQ ground state}
$\mathbf d_A=(1,0,0),\ \mathbf d_B=(0,1,0)$ give
$\langle Q_A^{x^2-y^2}\rangle=-1,\ \langle Q_B^{x^2-y^2}\rangle=+1$ (staggered
$(-1)^j$). At the SU(3) point the nearest-neighbor bond energy
$\propto|\mathbf d_i\!\cdot\!\bar{\mathbf d}_j|^2$ is minimized ($=0$ for orthogonal
directors) while the second-neighbor (same-sublattice, parallel) overlap is maximized
($=1$), matching the paper's stated minimization/maximization requirement.
\textbf{PASS.}

\subsection*{C4 --- Gell-Mann rotations and marginality structure}
The four generators (Eq.~S25) are Hermitian $\Rightarrow U(\phi)=e^{i\sum\lambda_p\phi_p}$
is unitary (residual $9\times10^{-16}$). For $200$ random rotations at the SU(3) point,
the two-director energy is invariant to $\sim10^{-15}$: the SU(3) rotations are exact flat
directions (Goldstone modes), the geometric substrate for the ``exactly marginal Kondo
coupling'' claim. \textbf{PASS} (structural; the full RG scaling
$S_K/S_c\propto\sqrt{\Lambda/K_F}$ is analytic and reproduced in the report text, not
independently re-derived numerically).

\subsection*{C5 --- SU(3) multiplet structure of the two-site spectrum}
At $J=K$ the two-site Hamiltonian has exactly two distinct eigenvalues with
degeneracies $\{6,3\}$ --- the $\mathbf 3\otimes\mathbf 3=\mathbf 6\oplus\bar{\mathbf 3}$
decomposition of SU(3). Away from $J=K$ ($J{=}1,K{=}1.2$) the levels split (three
distinct values), confirming the lowering SU(3)$\to$SU(2). \textbf{PASS.}

\subsection*{C6 --- Finite-cluster ED of the $(\pi,\pi)$ AFQ (DMRG surrogate)}
Exact diagonalization of the spin-1 BLBQ model at $J_1{=}1,K_1{=}1.2,J_2{=}0,K_2{=}-0.3$
on periodic $2\times2$ and $2\times4$ clusters (sparse Lanczos, cross-checked by dense
\texttt{eigh}): the quadrupolar structure factor $m_Q^2(\mathbf q)$ peaks at
$\mathbf q=(\pi,\pi)$ and dominates the dipolar $m_S^2$, whose weak peaks sit at
$(\pi,0)/(0,\pi)$ --- exactly the pattern of Figs.~1 and S1.

\begin{table}[h]\centering
\begin{tabular}{lccccc}
\toprule
cluster & $m_Q^2(\pi,\pi)$ & $m_S^2(\pi,\pi)$ & $m_S^2(\pi,0)$ & quad peak & spin peak\\
\midrule
$2\times2$ & $1.976$ & $0.384$ & $0.808$ & $(\pi,\pi)$ & $(\pi,0)$\\
$2\times4$ & $1.256$ & $0.218$ & $0.487$ & $(\pi,\pi)$ & $(\pi,0)$\\
\bottomrule
\end{tabular}
\caption{ED structure factors; dominant $m_Q^2$ at $(\pi,\pi)$, weak $m_S^2$ at
$(\pi,0)/(0,\pi)$, matching the paper qualitatively. Magnitudes are finite-size and are
\emph{not} claimed to match the thermodynamic-limit DMRG order parameters.}
\end{table}
\textbf{PASS} (qualitative; magnitudes out of scope for a laptop cluster).

\section{Verdict}
\begin{center}
\fbox{\parbox{0.9\textwidth}{\centering
\textbf{Verdict: REPLICATED (core algebra + qualitative numerics).}\\[2pt]
\textbf{Coverage: 7/10} --- all analytic/algebraic claims and the qualitative DMRG
structure-factor result reproduced; full DMRG magnitudes and the field-theoretic RG
derivation not independently reproduced (out of scope).\\[2pt]
\textbf{Agreement: 9/10} --- every in-scope quantitative check matches to machine
precision or with the correct qualitative peak structure.}}
\end{center}

All 12 algebraic sub-checks pass to $\lesssim10^{-15}$; the ED surrogate reproduces the
$(\pi,\pi)$-AFQ dominance on both clusters. No contradictions with the paper were found
within scope.

\end{document}
